%% Manuscrito SICST - Enfoque 1 (LLM vs humano)
%% Compilado con la plantilla oficial Springer Nature (sn-jnl.cls)
%% Descargada de: https://github.com/godkingjay/springer-nature-latex-template
%% (espejo publico de la plantilla oficial de descarga libre de Springer Nature)

\documentclass[pdflatex,sn-basic]{sn-jnl}

\usepackage{graphicx}
\usepackage{multirow}
\usepackage{amsmath,amssymb,amsfonts}
\usepackage{booktabs}
\usepackage{xcolor}
\usepackage{textcomp}
\usepackage{manyfoot}
\usepackage{float}
\usepackage{enumitem}

\raggedbottom

\begin{document}

\title[Human vs. LLM Requirements: A Paired Study]{An Empirical Comparison of Human-Elicited and LLM-Generated Functional Requirements in a Physical Therapy Follow-Up System: A Paired-Analyst Study in Ecuador}

\author*[1]{\fnm{Kevin Germán} \sur{Contreras Chávez}}\email{kcontrerasc3@uteq.edu.ec}

\author[1]{\fnm{Angelo Paul} \sur{Zambrano Moya}}\email{azambranom18@uteq.edu.ec}

\affil*[1]{\orgdiv{Software Engineering Program}, \orgname{Universidad Técnica Estatal de Quevedo (UTEQ)}, \orgaddress{\country{Ecuador}}}

\abstract{Digital support for home-based physical therapy depends on requirements that faithfully capture the needs of patients, caregivers, and physiotherapists. Large Language Models (LLMs) are increasingly proposed as assistants for requirements elicitation, but their comparative quality against human-elicited requirements remains underexplored in real-world, resource-constrained software engineering education settings. This paper reports a paired empirical study comparing the quality of functional requirements (FR) elicited by a human requirements engineering process against FR generated by an LLM (GPT-5.5-mini) from the same 18 anonymized interview transcripts, for the Intelligent System for Physical Therapy Control and Follow-up (SICST). The experimental protocol was registered retrospectively on the Open Science Framework (OSF, DOI 10.17605/OSF.IO/82Q76) on 2026-09-12 at 19:50 Ecuador time. The field interviews, generation of the LLM requirement set, and the four blind evaluations used in the analysis had already occurred before that registration. Consequently, the OSF record is not presented as a confirmatory preregistration and the inferential comparison is reported as exploratory. Thirty-three human-elicited and thirty-three LLM-generated FR (66 total) were anonymized and blind-rated by four raters across five quality dimensions (completeness, absence of ambiguity, verifiability, source correctness, and internal consistency) on a 1--5 Likert scale. Inter-rater agreement was quantified with Fleiss' kappa. Eleven requirement pairs with strict thematic correspondence were compared using paired $t$-tests, with Holm--Bonferroni correction for multiple comparisons. Inter-rater agreement was low to negligible across dimensions (Fleiss' $\kappa$ ranging from $-0.013$ to $0.167$). No dimension showed a statistically significant difference between human- and LLM-generated requirements after correction (all adjusted $p > 0.05$). These findings are interpreted as exploratory and inconclusive under the small paired sample, while highlighting the need for larger samples and more reliable evaluation protocols in future replications.}

\keywords{Requirements elicitation, Physiotherapy informatics, Large Language Models, Empirical study, Patient tracking}

\maketitle

\section{Introduction}

Physical therapy frequently extends beyond the clinical session. Patients may be expected to remember exercises, execute movements correctly, respect repetition and rest instructions, report pain or fatigue, and communicate problems between appointments. Telerehabilitation and digital rehabilitation systems can facilitate continuity of care, but their effectiveness and acceptability depend on how well technological interventions represent clinical and user needs \citep{munoztomas2023telerehabilitation,susomarti2021effectiveness,gava2022shoulder}. These activities therefore create a requirements problem before they create an algorithmic problem: a digital system must represent the needs, constraints, responsibilities, and safety expectations of patients, caregivers, and professionals in a way that can be verified and traced.

Requirements engineering (RE) provides the methods needed to elicit, analyze, document, validate, prioritize, and manage stakeholder expectations \citep{pohl2010requirements,iso29148_2018,sommerville2016software,pressman2020software}. Requirement quality depends not only on the final wording but also on clarity, verifiability, consistency, traceability, and the relationship between a requirement and its evidentiary source \citep{grosser2024templates,knauss2018clarification}. Understanding whether Large Language Models can support this elicitation step, and at what quality cost, matters directly to small teams that lack dedicated requirements analysts.

Recent work has begun to examine the use of generative AI for requirements engineering tasks, including elicitation, classification, refinement, and quality analysis \citep{cheng2026genai}. However, empirical comparisons based on real interview evidence remain limited, particularly in health-adjacent academic projects where privacy, traceability, human supervision, and methodological transparency are important.

SICST (\textit{Sistema Inteligente de Control y Seguimiento de Terapia Física}) was developed as an academic case study in Requirements Engineering at Universidad Técnica Estatal de Quevedo. The project addresses follow-up of physical therapy routines by separating patient and physiotherapist interactions while maintaining a common traceability chain from field evidence to requirements, design elements, and prototype behavior. Computer-vision approaches have also been investigated broadly for rehabilitation and movement assessment, illustrating the need to define measurable technical, ethical, and human-supervision requirements before such functionality is incorporated into health-adjacent systems \citep{debnath2022rehabilitation,oecd2019ai}.

This paper reports the terminal empirical component of the project: a comparison of requirements quality between the team's own elicitation process and an LLM prompted with the same underlying field evidence.

This study contributes: (1) a documented blind comparison protocol for human- and LLM-generated functional requirements in a real, non-simulated domain, accompanied by a retrospectively registered OSF protocol whose actual timing is reported transparently; (2) an exploratory empirical estimate of the quality differences between both sources across five RE quality dimensions; and (3) a transparent report of inter-rater reliability and provenance limitations when four non-expert raters evaluate requirements quality on a Likert scale.

\begin{description}[leftmargin=1.3cm,labelwidth=1.0cm]
\item[RQ1.] Do quality dimensions (completeness, absence of ambiguity, verifiability, source correctness, internal consistency) differ between functional requirements elicited by a human requirements engineering process and those generated by a Large Language Model, when both are derived from the same anonymized interview transcriptions?
\end{description}

\noindent\textbf{H0.} There are no statistically significant differences in the quality of functional requirements generated through a human process versus those generated through a language model.

\noindent\textbf{H1.} There are statistically significant differences in the quality of functional requirements generated through a human process versus those generated through a language model.

\section{Background and Related Work}

\subsection{LLMs for requirements elicitation}

Generative language models have been proposed for several RE sub-tasks, including requirements extraction from unstructured text, ambiguity detection, refinement, and requirements quality assessment \citep{cheng2026genai}. Their ability to generate fluent technical language makes them potentially useful as drafting assistants, but fluent wording does not by itself demonstrate correctness, traceability, or completeness.

Requirements engineering literature emphasizes that elicited statements must remain connected to stakeholder needs and must evolve through clarification and validation rather than being treated as isolated text-generation outputs \citep{pohl2010requirements,knauss2018clarification}. Requirement templates can improve consistency and testability, although their appropriateness and expressiveness depend on the context in which they are used \citep{grosser2024templates}.

Evidence on LLM-assisted elicitation quality remains comparatively limited. Many evaluations rely on synthetic, benchmark, or publicly available corpora rather than transcripts from a live field study. This motivates transparent empirical comparisons in which both human and LLM requirement sets originate from the same evidence.

\subsection{Requirements quality dimensions}

Completeness, absence of ambiguity, verifiability, traceability to source, and internal consistency are long-standing quality concerns for requirements statements \citep{iso29148_2018,sommerville2016software,pressman2020software}. ISO/IEC/IEEE 29148 emphasizes characteristics of well-formed requirements and the importance of verification and traceability \citep{iso29148_2018}.

We adopt five dimensions adapted to allow blind comparison between two requirement sources drawn from the same evidentiary base, rather than applying the dimensions as an absolute compliance checklist. The evaluation therefore reflects perceived quality under the instrument and rater population used in this study.

\subsection{Physical therapy follow-up and telerehabilitation}

Digital follow-up and telerehabilitation can complement face-to-face physical therapy and provide mechanisms for remote therapeutic exercise and monitoring \citep{munoztomas2023telerehabilitation,susomarti2021effectiveness}. Evidence across rehabilitation contexts also highlights the importance of appropriately designed remote interventions rather than assuming technology itself guarantees therapeutic benefit \citep{gava2022shoulder}.

Digital physical-therapy systems introduce constraints beyond generic software functionality, including exercise adherence, pain and fatigue reporting, patient privacy, caregiver involvement, and professional oversight. Computer-vision approaches for rehabilitation further introduce performance, fairness, explainability, and supervision concerns \citep{debnath2022rehabilitation,chazette2020explainability}. These considerations motivated the original field study from which both requirement sets were derived.

\subsection{Inter-rater agreement in requirements evaluation}

Fleiss' kappa is a standard measure of agreement among three or more raters using categorical assessments, while Cohen's kappa is commonly used for agreement between two raters \citep{fleiss1971measuring,cohen1960coefficient}. These measures are useful because percentage agreement alone does not account for agreement expected by chance.

Low agreement in subjective quality ratings is a significant risk in small-panel, non-expert evaluation designs. The present study therefore reports agreement directly rather than attempting to remove or conceal disagreement among evaluators.

\subsection{Empirical-study transparency}

Empirical software-engineering research requires transparent reporting of study design, evidence, analysis decisions, and limitations \citep{runeson2009guidelines,wohlin2012experimentation}. In qualitative work, documentation of sampling, coding, and thematic saturation also helps readers judge the strength and scope of conclusions \citep{guest2006interviews,hennink2017codesaturation}.

In this study, transparency also requires distinguishing clearly between a genuine preregistration and a protocol that was registered retrospectively after experimental artifacts already existed. That distinction is important when interpreting whether statistical analyses should be considered confirmatory or exploratory.

\section{Methodology}

\subsection{Design}

This is a non-experimental, paired comparative study. No causal relationship between requirement source and quality is inferred; the study only compares the perceived quality of two requirement sets drawn from the same source evidence.

The experimental protocol was registered on OSF (DOI \href{https://doi.org/10.17605/OSF.IO/82Q76}{10.17605/OSF.IO/82Q76}) at 2026-09-13T00:50:38Z, corresponding to 2026-09-12 19:50:38 in continental Ecuador.

The chronology reconstructed from the repository evidence is:

\begin{enumerate}[leftmargin=1.2cm]
\item field interviews and requirements elicitation;
\item generation of the LLM requirement set;
\item completion of the four blind-rating worksheets;
\item OSF registration of the experimental protocol;
\item subsequent statistical analysis and documentation.
\end{enumerate}

Therefore, the field interview corpus, the LLM-generated requirement set, and the four blind-rating spreadsheets used in the analysis already existed before the OSF registration timestamp.

The OSF record is consequently documented as a \textbf{retrospective protocol registration}, not as a confirmatory preregistration preceding the generation of experimental data. All inferential comparisons reported in this manuscript are interpreted as exploratory.

\subsection{Population and sampling}

The unit of analysis is the individual functional requirement (FR). Two comparison groups were constructed from the same evidentiary source:

\begin{itemize}[leftmargin=1.2cm]
\item \textbf{Human-elicited group} ($n = 33$): FR produced through the team's standard requirements engineering process involving interviews, thematic analysis, and consolidation into the project's ERS/SRS.
\item \textbf{LLM-generated group} ($n = 33$): FR produced by prompting GPT-5.5-mini with the same 18 anonymized interview transcripts used for human elicitation, under an instruction set requiring each requirement to be grounded exclusively in the supplied transcripts (see Section~\ref{sec:llm-prompt}).
\end{itemize}

Both sets were combined into a single list of 66 anonymized items and presented to four raters who were blind to requirement origin (Section~\ref{sec:blinding}).

\subsection{LLM requirements generation}
\label{sec:llm-prompt}

The LLM (GPT-5.5-mini, accessed via ChatGPT) was given the full text of the 18 anonymized interview transcripts and an explicit rule set requiring: (a) grounding every requirement strictly in the supplied transcripts, with no external knowledge or invented functionality; (b) a fixed sentence template (\textit{``The system shall [verifiable action] [object/condition]''}); (c) specificity, verifiability, and traceability to a concrete transcript fragment; (d) no duplicate or overlapping requirements; (e) sequential numbering (RF-01 to RF-33); and (f) output restricted to a CSV with columns \texttt{id\_requisito, descripcion}, with no additional prose.

A second turn asked the model to confirm its version and confirm that no external information was used. The prompt, both documented turns, the model confirmation, and the resulting CSV are versioned in \texttt{06\_Experimento/prompts\_llm/}.

The exact clock time of the LLM generation was not preserved. Repository evidence nevertheless constrains it to before the blind-evaluation spreadsheets were saved on 2026-09-12 between approximately 13:11 and 13:16 Ecuador time. The previously documented claim that LLM generation occurred on 2026-09-13 after OSF registration was corrected during the final review.

\subsection{Blinding and evaluation instrument}
\label{sec:blinding}

The 66 FR (33 human and 33 LLM) were normalized to a common format, shuffled, and assigned anonymous item identifiers with no indication of source. Four raters, blind to the origin of each item, rated every requirement on a 1--5 Likert scale across five dimensions: completeness, absence of ambiguity, verifiability, source correctness, and internal consistency.

The worksheet presented to the raters contained only anonymous ITEM identifiers and requirement text, with no source label. After the evaluations had been completed, the mapping ITEM $\rightarrow$ real requirement identifier $\rightarrow$ source (human/LLM) was reconstructed from the versioned human traceability matrix and the LLM requirement set and exported as \texttt{06\_Experimento/resultados/mapa\_origen\_items\_reconstruido.csv} for reproducible analysis.

The four evaluations were performed sequentially during the same meeting using Angelo Zambrano's computer. Each evaluator completed his or her own worksheet. After an evaluator finished, Angelo checked only that the rating fields had been completed and saved the file before continuing with the next evaluator.

The procedure was intended to avoid raters consulting one another's recorded scores. However, because the evaluations occurred in the same session and on a single computer, the provenance evidence is weaker than it would have been if each evaluator had completed and returned a separate file from an independent device. The common last-editor metadata in the four spreadsheets is therefore consistent with the final saving of the files on the organizer's computer, but that metadata alone does not establish who entered each individual rating.

\subsection{Thematic pairing}

To support a paired statistical comparison, human and LLM requirements addressing the same underlying interview theme were matched by two reviewers into a traceability matrix (\texttt{matriz\_trazabilidad\_tema\_RF.csv}). Eleven pairs achieved strict, unambiguous thematic correspondence (type \texttt{pareado}); these are the pairs used in the main exploratory paired analysis reported in Section~\ref{sec:results}.

The remaining items, including partial matches, weak matches, or items with no thematic counterpart in the other set, are retained in the traceability matrix for transparency but excluded from the paired test because they do not support a valid 1-to-1 comparison.

\subsection{Statistical analysis plan}

For each of the five quality dimensions and for the aggregate score, the analysis: (i) computed Fleiss' kappa across all four raters and 66 items to quantify inter-rater agreement; (ii) tested normality of the 11 within-pair differences (human $-$ LLM) with the Shapiro--Wilk test; (iii) applied a paired $t$-test when normality was not rejected, or a Wilcoxon signed-rank test otherwise; and (iv) corrected the five dimension-level $p$-values for multiple comparisons using the Holm--Bonferroni procedure.

A 95\% bootstrap confidence interval with 10{,}000 resamples was computed for the mean paired difference. All analyses were executed by a versioned script (\texttt{06\_Experimento/scripts\_analisis/analisis\_experimento\_llm\_humano.py}) that reproduces the reported quantitative outputs directly from the raw blind-rating spreadsheets.

\subsection{Retrospective registration and protocol deviations}

The OSF registration occurred after the field interviews, generation of the LLM requirement set, and collection of the four blind-rating worksheets. It therefore cannot be treated as a confirmatory preregistration of those activities.

The exact statistical test per dimension (paired $t$ vs.\ Wilcoxon) and the Holm--Bonferroni correction were operationalized after registration because the OSF protocol described the general comparison strategy without fixing the individual statistical test names.

Accordingly, the study does not use the OSF timestamp to claim that the LLM generation, blind evaluations, or resulting data collection were prospectively preregistered. These timing and protocol deviations are reported transparently, and all inferential results in this manuscript are interpreted as exploratory rather than confirmatory.

\section{Results}
\label{sec:results}

\subsection{Inter-rater agreement}

Agreement among the four blind raters was low to negligible across all five dimensions (Table~\ref{tab:kappa}). This indicates that the raters frequently disagreed on how to score individual requirements on the 1--5 scale, independently of requirement source.

\begin{table}[h]
\centering
\caption{Fleiss' kappa by quality dimension (all 66 items, 4 raters).}
\label{tab:kappa}
\begin{tabular}{lc}
\toprule
Dimension & Fleiss' $\kappa$ \\
\midrule
Completeness & 0.132 \\
Absence of ambiguity & 0.167 \\
Verifiability & $-0.013$ \\
Source correctness & $-0.011$ \\
Internal consistency & 0.046 \\
\botrule
\end{tabular}
\end{table}

\subsection{Supporting field-study descriptives and qualitative reliability}

The reproducible field-data package under \texttt{07\_Datos/resultados/} contains 79 questionnaire responses: 62 patients or former physical-therapy patients, 14 family caregivers, and 3 physiotherapists.

Among the 76 responses to the camera-acceptance item, 43 (56.6\%) selected \textit{yes, but only with prior authorization} and 21 (27.6\%) selected \textit{yes}. Among the 76 responses to the progress-visualization item, 64 (84.2\%) answered yes.

The independent qualitative double-coding check compared 36 coding decisions, with 33 agreements and 3 disagreements (91.67\% raw agreement). Cohen's $\kappa$ was 0.7187 with a 95\% bootstrap interval of [0.3077, 1.0000], based on 10{,}000 resamples with seed 401. This statistic describes agreement between the two qualitative coders and is distinct from the Fleiss' $\kappa$ values above, which describe agreement among the four blind requirement raters.

The thematic-saturation analysis accumulated 17 themes across the 18 canonical elicitation sessions. The last new theme appeared at session 9 (EFT-01); sessions 10--18 contributed no additional themes. Saturation is reported descriptively rather than treated as proof that every possible stakeholder need had been captured, consistent with established distinctions between code saturation and deeper meaning saturation in qualitative research \citep{guest2006interviews,hennink2017codesaturation}.

\subsection{Paired comparison (RQ1)}

Table~\ref{tab:pareado} reports the paired analysis over the 11 thematically-matched requirement pairs. All five within-pair difference distributions did not reject normality (Shapiro--Wilk $p > 0.05$), so a paired $t$-test was used throughout.

After Holm--Bonferroni correction for the five dimension-level comparisons, no dimension reached statistical significance.

\begin{table}[h]
\centering
\caption{Paired comparison, human-elicited vs.\ LLM-generated FR ($n=11$ pairs). Mean difference is human minus LLM.}
\label{tab:pareado}
\small
\begin{tabular}{lccccc}
\toprule
Dimension & $t(10)$ & Mean diff.\ & 95\% CI & $p$ (raw) & $p$ (Holm) \\
\midrule
Completeness & 0.714 & 0.091 & [$-0.136$, 0.341] & 0.492 & 1.000 \\
Absence of ambiguity & $-0.468$ & $-0.091$ & [$-0.455$, 0.250] & 0.650 & 1.000 \\
Verifiability & 1.000 & 0.136 & [$-0.114$, 0.386] & 0.341 & 1.000 \\
Source correctness & 1.550 & 0.159 & [$-0.045$, 0.341] & 0.152 & 0.761 \\
Internal consistency & 0.938 & 0.091 & [$-0.091$, 0.273] & 0.371 & 1.000 \\
\botrule
\end{tabular}
\end{table}

Figure~\ref{fig:pareado} shows the paired global scores, calculated as the average of the five dimensions, for each of the 11 matched themes.

\begin{figure}[h]
\centering
\includegraphics[width=0.75\linewidth]{figuras/figura_comparacion_pareada.png}
\caption{Paired global quality score by thematic pair, human-elicited vs.\ LLM-generated requirements ($n=11$).}
\label{fig:pareado}
\end{figure}

We cannot reject H0 for any of the five quality dimensions. The complementary global paired score also did not show a statistically significant difference (mean human-minus-LLM difference 0.0773, 95\% bootstrap CI [$-0.0455$, 0.1955], $t(10)=1.193$, $p=0.261$).

Within the limits of this small-sample exploratory design, the observed ratings do not provide evidence of a difference in perceived quality between the two sources. This absence of statistically significant evidence must not be interpreted as evidence of equivalence.

\section{Discussion}

\subsection{Implications for research}

The absence of a detectable quality gap, combined with very low inter-rater agreement, is an informative methodological result. It shows that under this particular instrument, sample, and group of raters, the study did not obtain stable evidence of a difference between the two requirement sources.

LLM-generated technical text can resemble the surface form of conventional requirements documentation, while correctness and traceability continue to require explicit verification against source evidence \citep{cheng2026genai,iso29148_2018}. Future comparisons should therefore combine textual quality assessment with stronger traceability checks and better-calibrated evaluation instruments.

\subsection{Implications for practice}

For small teams without dedicated requirements specialists, an LLM can serve as a drafting aid when it is constrained to documented source material. However, this study does not provide evidence sufficient to recommend replacing human elicitation with automated generation.

Human review, stakeholder validation, traceability to source evidence, and explicit verification remain necessary parts of requirements engineering \citep{pohl2010requirements,sommerville2016software,pressman2020software}.

\subsection{A surprising or counterintuitive finding}

The negative and near-zero Fleiss' kappa values for verifiability and source correctness are notable: on those dimensions, the four raters showed little agreement under the instrument employed.

Rather than treating low agreement as evidence about either requirement source alone, it should also motivate improvement of the evaluation procedure itself. Future replications could use behaviorally anchored examples for each Likert point, evaluator training, calibration exercises, or a smaller panel with greater RE expertise.

We restrict all conclusions to this specific, small-sample, Ecuadorian, physical-therapy domain study. No claim is made that LLMs equal, exceed, or replace human analysts in general software-engineering settings.

\section{Threats to Validity}

Following common software-engineering empirical-study practice \citep{runeson2009guidelines,wohlin2012experimentation}, threats are organized into construct, internal, external, and conclusion validity.

\subsection{Construct validity}

The five quality dimensions operationalize subjective judgments on a 1--5 scale without behaviorally anchored examples for every rating level. This may have contributed to the low inter-rater agreement observed.

The main exploratory paired test uses only the 11 strictly matched pairs. The broader, weaker, or unmatched items---44 of the 66 requirements---are not part of the paired statistical claim. Therefore, results for the 11 pairs should not be interpreted as an evaluation of every requirement in both sets.

The comparison also evaluates perceived requirement quality under the selected dimensions and instrument; it does not directly measure downstream software correctness, clinical effectiveness, patient outcomes, or the overall competence of either a human analyst or an LLM.

\subsection{Internal validity}

Several internal-validity limitations are relevant.

First, the LLM-generated requirement set and the four blind-rating worksheets already existed before the OSF registration timestamp. The OSF record is therefore retrospective with respect to the experiment. It cannot establish that generation choices, rating procedures, or hypotheses were fixed prospectively before experimental data existed. For that reason, the inferential comparison is treated as exploratory rather than as a preregistered confirmatory analysis.

Second, the four evaluations were performed sequentially during the same meeting using the same computer. According to the documented procedure, each evaluator completed his or her own worksheet and the organizer checked only that the rating fields were complete before saving the file. The arrangement was intended to prevent evaluators from consulting previously entered ratings, but the study does not have separate-device or independently transmitted files that would provide stronger provenance evidence.

The internal metadata of the four spreadsheets consequently identifies the same last editor and records saves within a short interval. These metadata are consistent with the documented same-device saving procedure, but they do not independently prove which person entered each individual rating. This limitation is explicitly retained rather than inferred away.

Third, the raters were blind to requirement origin but were not trained requirements engineers. Their limited expertise and the absence of a calibration phase may explain part of the low Fleiss' kappa values.

Fourth, the precise statistical test-selection rules and Holm--Bonferroni procedure were operationalized after the OSF registration. Although the resulting analysis is reproducible from the versioned script, the chronology prevents those decisions from being treated as part of a prospective confirmatory protocol.

\subsection{External validity}

The work concerns one academic project, one LLM, one prompting strategy, one Spanish-language interview corpus, one physical-therapy application domain, and a specific Ecuadorian educational context.

Results therefore should not be generalized directly to other LLMs, languages, prompting strategies, organizations, domains, professional requirements engineers, or clinical populations without additional replication.

The health-adjacent context also introduces privacy, supervision, and explainability concerns that may not be present to the same degree in other software domains \citep{chazette2020explainability,oecd2019ai,ecuador2021lopdp}.

\subsection{Conclusion validity}

Only 11 requirement pairs satisfied the strict thematic matching criterion. Statistical power to detect small-to-moderate effects is consequently limited.

The failure to reject H0 cannot be interpreted as evidence that both sources are equivalent. It indicates only that this study, with this sample, instrument, rater panel, and analysis, did not detect statistically significant differences after correction for multiple comparisons.

Low inter-rater agreement introduces additional uncertainty because the measured scores themselves vary substantially across evaluators. Accordingly, the study's strongest conclusion is methodological and exploratory rather than a claim of human--LLM equivalence.

\section{Ethics, Privacy, and Data Governance}

The underlying interview transcripts used to generate both requirement sets were collected under the project's documented informed-consent procedures and stored in anonymized form in the research materials, while direct identifiers and signed consent evidence are restricted to an encrypted, access-controlled container.

The blind requirement-evaluation worksheet contained anonymous ITEM identifiers and requirement text rather than human participant identifiers. The source mapping used for analysis was reconstructed after the ratings were collected and is published as a reproducibility artifact.

Privacy and governance considerations are relevant because the project concerns health-adjacent information. The project therefore separates publicly reproducible research artifacts from direct identifiers and restricted evidence, consistent with the need to minimize unnecessary disclosure of personal information \citep{ecuador2021lopdp}.

\section{Reproducibility and Open Materials}

The experimental package is versioned in the project repository under \texttt{06\_Experimento/}. It contains the protocol (\texttt{protocolo.pdf}), OSF registration evidence, deviations documentation, exact LLM prompt record, raw blind-rating spreadsheets, LLM-generated requirements, thematic pairing matrix, analysis script, and generated results.

The OSF record must be interpreted according to its actual chronology: it is a retrospective registration of the experimental protocol and not evidence that the LLM generation or blind evaluations were prospectively preregistered.

The analysis script under \texttt{06\_Experimento/scripts\_analisis/} reproduces the reported comparison outputs from the deposited raw rating data. Supporting field-data analyses and their derived results are versioned under \texttt{07\_Datos/}.

\begin{center}
Dataset DOI: \url{https://doi.org/10.5281/zenodo.22315298}
\quad\textbar\quad
OSF retrospective protocol registration:
\url{https://doi.org/10.17605/OSF.IO/82Q76}
\end{center}

\section{Conclusion}

This paper reported a blind comparison of the quality of functional requirements elicited by a human process versus generated by an LLM (GPT-5.5-mini) for a physical-therapy follow-up system, using the same 18 anonymized interview transcripts as the evidentiary source.

The experimental protocol was registered retrospectively on OSF. The field interviews, documented LLM-generation phase, and four blind evaluations all predated the OSF registration timestamp. The OSF record is therefore not presented as a prospective confirmatory preregistration, and the statistical comparison is interpreted as exploratory.

Inter-rater agreement among the four blind raters was low to negligible across all five evaluated quality dimensions. After correcting for multiple comparisons, no dimension showed a statistically significant difference between human-elicited and LLM-generated requirements (all adjusted $p > 0.05$); H0 could not be rejected.

These results, obtained under a small paired sample ($n=11$), should be read as an inconclusive empirical finding rather than as evidence that LLMs match human requirements engineers. Future work should increase the number of thematically matched pairs, use independently collected evaluator files, train and calibrate raters, establish prospective analysis protocols before collecting experimental ratings, and replicate the comparison across additional models, domains, and participant populations.

\section*{Declarations}

\textbf{Funding.}
This academic project reports no external research funding in the materials used for this manuscript.

\textbf{Conflict of interest.}
The authors declare no conflict of interest relevant to this manuscript.

\textbf{Author contributions.}
Angelo Paul Zambrano Moya conducted the OSF protocol registration, executed the documented LLM requirement-generation process, coordinated the blind-rating activity, performed the statistical analysis, and participated in methodological documentation and repository closure. Kevin Germán Contreras Chávez participated as a remote co-interviewer during field evidence collection and contributed during the final repository-closure phase through evidence organization, preparation of 22 retrospective syntheses from existing transcripts, correspondence, EXIF inventory work, pre-delivery verification, individual evidence, and the two jointly recorded work sessions. Both authors reviewed the manuscript and checked its claims against the versioned evidence and reproducible project artifacts.

\textbf{Data availability.}
The reproducible dataset is available at Zenodo, DOI:
\url{https://doi.org/10.5281/zenodo.22315298}.
The retrospectively registered experimental protocol is available at OSF, DOI:
\url{https://doi.org/10.17605/OSF.IO/82Q76}.
Public materials exclude direct identifiers and signed consent documents.

\textbf{Code availability.}
The project source repository, including the analysis scripts and versioned research artifacts, is publicly available at:
\url{https://github.com/gleiston-guerrero/SICST_ISR401}.

\textbf{Use of AI tools.}
An LLM (GPT-5.5-mini) was used to generate one of the two compared requirement sets as the object of study itself; this use is documented in \texttt{06\_Experimento/prompts\_llm/}. AI assistance was also used during repository organization, technical-writing review, data-processing script development, and manuscript drafting from project data and results. No participant, interview response, questionnaire observation, rating, or empirical result is claimed to have been generated as substitute evidence by this assistance.

\bibliography{references}

\end{document}